\section{静电场 III}

\subsection{电势}

\subsubsection{静电场的环路定理}

\begin{theorem}
	在静电场中，沿环路电场力所做的功为 $0$。
\end{theorem}

\subsubsection{电势}

在静电场中的两点 $A, B$，电势差 $\varphi_A - \varphi_B$ 为单位正电荷 $e$ 从 $A \to B$ 所受电场力做的功。

取定一个势能零点 $O$，那么可以定义静电场的电势：

\begin{definition}[电势]
	点 $P$ 的电势定义为
	$$
	\int_P^{O} \vect{E} \vdot \dd{\vect{l}}
	$$
\end{definition}

电势同样满足叠加原理。

\subsection{静电场的基本微分方程}

\subsubsection{矢量分析}

\begin{definition}[散度和旋度]
	对于一个矢量场函数 $\vect{F}(x, y, z) = P \vect{i} + Q \vect{j} + R \vect{k}$，定义
	\begin{itemize}
		\item \textbf{散度}为 $\displaystyle \div \vect{F} = \pdv{P}{x} + \pdv{Q}{y} + \pdv{R}{z}$；
		\item \textbf{旋度}为 $\displaystyle \curl \vect{F} = \begin{vmatrix}
			\vect{i} & \vect{j} & \vect{k} \\
			\pdv{x} & \pdv{y} & \pdv{z} \\
			P & Q & R 
		\end{vmatrix} = \ab(\pdv{R}{y} - \pdv{Q}{z}) \vect{i} + \ab(\pdv{P}{z} - \pdv{R}{x}) \vect{j} + \ab(\pdv{Q}{x} - \pdv{P}{y}) \vect{k}
		$；
	\end{itemize}
\end{definition}


\subsubsection{静电场高斯定理和环路定理的微分形式}

% $$
% \begin{dcases}
	
% \end{dcases}
% $$

$$
\begin{dcases}
	\div \vect{E} = \frac{\rho}{\varepsilon_0} \\
	\curl \vect{E} = \vect{0}
\end{dcases}
$$

这个方程称为\textbf{静电场基本方程的微分形式}。